\chapter*{Start here: from zero}
\addcontentsline{toc}{chapter}{Start here: from zero}
\markboth{Start here}{Start here}

This book is about answering one question honestly: \emph{does this trading idea
really work, or does it only look like it does?} You do not need to know finance.
You need comfort with percentages, a little statistics, and some Python. This
chapter covers the finance vocabulary and the statistics you need for Chapter~1.

\section*{What a market is}

A \emph{market} is a place where buyers and sellers agree on prices. An
\emph{asset} is whatever is traded: a share of a company, a currency, a
cryptocurrency. A \emph{price} is the amount paid per unit at one moment. A
\emph{trade} swaps cash for the asset or the other way around; a \emph{strategy}
is a rule that says when to trade. \emph{Quantitative research} tests such rules
on data, with statistics, before any money is risked.

This book is a research-methods text. Nothing in it is financial advice, and the
examples are about measurement, not about what to buy.

\section*{Returns: the one idea you need first}

If an asset costs $100$ today and $102$ tomorrow, its \emph{return} is the
percentage change, $102/100-1=0.02$, or $2\%$. If it then falls to $99$, the next
return is $99/102-1=-0.0294$, a fall of $2.94\%$. Two things to notice. First,
the up-then-down does not cancel: a $+2\%$ move followed by a fall of $2.94\%$
returns to roughly $99$, not $100$. Second, we will often use the
\emph{logarithmic return} $\ln(P_{t}/P_{t-1})$ instead, because log returns
\emph{add} over time. For small moves the two are nearly equal ($\ln 1.02=0.0198$
versus $0.02$); for large moves they differ ($\ln 1.1=0.0953$ versus $0.10$).
Chapter~1 explains exactly when the difference matters.

\section*{Three statistics you will use constantly}

\begin{description}
\item[Mean] The average. For the returns $+2\%,\,-2.94\%,\,+1\%$ it is
      $0.0196\%$ per period ($0.000196$).
\item[Standard deviation] How far values typically sit from the mean: the
      spread, or in finance the \emph{volatility}. For those three returns it is
      about $2.6\%$ (sample standard deviation $0.0261$).
\item[Correlation] A number between $-1$ and $1$ measuring whether two series move
      together. It is used again in Chapters~14--16.
\end{description}
\textbf{Book~3, Chapters~12--13} derive all three carefully, and each chapter here
recaps the idea at the point it is first used.

\section*{Python}

All code in this book is Python with NumPy and pandas. If you can read a loop and
a function, that is enough; the machine-learning book (Book~4) opens with a Python
primer you can read first.

\section*{How the books of this series fit together}

\begin{center}
\begin{tabular}{P{2.4cm}P{9.0cm}}
\toprule
\thead{Book} & \thead{Connection}\\
\midrule
3 Linear algebra & Statistics (Ch.~12--13), covariance matrices (Ch.~9--10) and optimisation (Ch.~14) are the maths under Chapters~13--16 here.\\
4 Machine learning & Shares the honest-evaluation discipline: train/test splits and leakage (Book~4, Ch.~3) are the same ideas as walk-forward validation (Chapter~11 here).\\
1 Competitive programming & The coding fluency used for the backtest harness (Chapter~8) and the exchange simulator (Chapter~3).\\
\bottomrule
\end{tabular}
\end{center}

\section*{How to use the boxes in each chapter}

Each chapter opens with \emph{After this chapter you can}. From Chapter~2, a
\emph{Before you start} box lists what the chapter relies on, with a recap. If a
number in the book surprises you, rerun the code that produced it; every number
is meant to be checkable.
